\begin{lemma}[Scheide single boundary crossing (20)]
Let $\alpha$ be present at every vertex of the enlarged set $X$.  In the
full $(\alpha_0,\alpha)$ incidence component through $y_0$, there is exactly
one individual edge identity in $\partial_G(X)$.
\end{lemma}
\begin{proof}
Put
\[
 X_0=\noderef{TreeSequenceVertexSet}(T),\qquad
 X=\noderef{ScheideEnlargedVertexSet}(R),\qquad
 y_0=C.y(0).
\]
Unpack the inherited \noderef{ScheideStructuralConfiguration} in the
reduced configuration \noderef{ScheideReducedConfiguration}.  Its final
field is a globally maximum triple.  By
\noderef{ScheideGlobalMaximumOrder}, $X_0$ is elementary and closed under
$\varphi$.  The \noderef{GlobalMaximumTriple} field also contains a
\noderef{KTriple}, so $\varphi$ is proper in the sense of
\noderef{IsProperEdgeColoring}.  The exact base-position equation and
\noderef{TreeSequenceVertexSet} put $y_0$ in $X_0$, hence in $X$.
The setup \noderef{ScheideDefectiveSetup} gives
\[
 \alpha_0\in\noderef{MissingColors}(G,\varphi,y_0).          \tag{1}
\]
The enlarged set is elementary by \noderef{ScheideEnlargedElementary}
and closed by \noderef{ScheideEnlargedClosed}.

By hypothesis, $\alpha$ belongs to
\noderef{PresentColors}$(G,\varphi,y_0)$, whereas (1) and the complementary
definition of \noderef{MissingColors} say that $\alpha_0$ does not.
Thus $\alpha\ne\alpha_0$.  Properness permits at most one incident edge
identity of either color at any vertex: two distinct such identities
would be adjacent by \noderef{AdjacentEdges}.  It follows from
\noderef{TwoColorIncidenceGraph} that the vertex-node $y_0$ has degree
one in its $(\alpha_0,\alpha)$ component.  Every edge-node in that
component has exactly its two endpoint vertex-neighbors: this is the
definition of the incidence graph together with \noderef{Incident} and
the distinct-end law in \noderef{LooplessMultigraph}.  All nodes have
degree at most two also by \noderef{TwoColorIncidenceDegreeAtMostTwo}.

Apply \noderef{TwoColorComponentSpanningPath} to this component.  Orient
the resulting simple spanning path $P$ from its degree-one node $y_0$.
The degree bounds show that the component has no incidence outside this
path: an extra incidence at an internal path node would give degree at
least three, while an extra incidence joining the two path ends would
give degree two at $y_0$.  Thus the path contains all nodes and all
incidences of the full set \noderef{TwoColorComponent} through $y_0$.
Its other end is a vertex-node, since every edge-node has degree two.
If that other end $w$ belonged to $X$, then $\alpha$ would be present at
$w$ by hypothesis.  Since $w$ has degree one, $\alpha_0$ would be missing
there.  The two distinct vertices $y_0,w\in X$ would then both miss
$\alpha_0$, contrary to elementariness and
\noderef{ColoredVertexSetElementary}.  Consequently the other end of
$P$ lies outside $X$.

Let $v$ be the last vertex of $X$ on $P$, let $g$ be the next edge-node,
and let $w$ be the vertex-node after $g$.  Then $v\in X$, $w\notin X$,
so $g\in\noderef{BoundaryEdges}(G,X)$.  The edge $g$ has color
$\alpha_0$ or $\alpha$ by the definition of the incidence graph.
Closedness of $X$, expanded through
\noderef{ColoredVertexSetClosed}, and (1) forbid an
$\alpha_0$-colored boundary identity.  Hence
\[
 \varphi(g)=\alpha.                                         \tag{2}
\]
We prove in three cases that the segment $y_0Pv$ lies in $X$.

First suppose that $v\in X_0$ and that $\alpha$ is not in
\noderef{DefectiveColors}$(G,\varphi,X_0)$.  The edge $g$ has one end in
$X_0$ and its other end outside $X\supseteq X_0$, so (2) makes it an
$\alpha$-colored boundary identity of $X_0$.  Original closedness and
the missing occurrence (1) forbid every $\alpha_0$-colored boundary
identity of $X_0$.  If $P$ left $X_0$ before reaching $v$, its first such
departure would therefore be an $\alpha$-colored boundary identity.
It is distinct from $g$ because $P$ is simple.  This would give two
members of \noderef{BoundaryColorEdges}$(G,\varphi,X_0,\alpha)$, which
is exactly defectiveness by \noderef{DefectiveColors}.  The contradiction
proves
\[
 y_0Pv\subseteq X_0\subseteq X.                              \tag{3}
\]

Next suppose that $v\in X_0$ and that $\alpha$ is defective for $X_0$.
Regard $\alpha$ as the defective-color index $d$, put
\[
 \gamma= C.\gamma(D.\gammaIndex(d)),\quad
 Z=X_0\setminus U_\gamma,\quad
 U_\alpha=\{(R d).u1,(R d).u2\},\quad X_\alpha=X_0\cup U_\alpha.
\]
Here $U_\gamma$ is the two-vertex block deleted in the exact replacement
data \noderef{ScheideReplacementData}.  The first conclusion of
\noderef{ScheideSymmetricDefectiveInternal} supplies an
$\alpha$-colored identity internal to $U_\gamma$.  Its two ends are the
two distinct members of $U_\gamma$ by looplessness.  Properness therefore
excludes any other $\alpha$-edge incident with either of those vertices.
Thus every $\alpha$-colored boundary identity of $X_0$ has its tree end
in $Z$.  The exact-three clause of the replacement data now says that
these identities are precisely
\[
 f^0_\alpha,\quad f^1_\alpha,\quad f^2_\alpha.               \tag{4}
\]
The latter two have their outside-$X_0$ ends in $U_\alpha$, so they are
internal to $X_\alpha$.  They are distinct and have color $\alpha$.
Their two ends in $U_\alpha$ are distinct, since otherwise the two named
edges would be distinct incident edges of the same color, contrary to
properness.  Thus each member of $U_\alpha$ is already incident with its
named $\alpha$-edge $f^1_\alpha$ or $f^2_\alpha$, and properness excludes
any further $\alpha$-edge incident with that member.  In particular, no
$\alpha$-colored boundary identity of $X_\alpha$ can leave from
$U_\alpha$.  The setup puts the outside end of
$f^0_\alpha$ outside $X_0$.  It differs from both members of $U_\alpha$:
otherwise $f^0_\alpha$ and one of the distinct identities
$f^1_\alpha,f^2_\alpha$ would be two incident $\alpha$-edges, contrary
to properness and \noderef{AdjacentEdges}.  Hence $f^0_\alpha$ leaves
$X_\alpha$, and (4) proves that it is the unique $\alpha$-colored
boundary identity of $X_\alpha$.

The set $X_\alpha$ is closed under $\varphi$ by
\noderef{ScheideReplacementUnionClosed}, and it contains $y_0$, so (1)
forbids every $\alpha_0$-colored boundary identity of $X_\alpha$.
The edge $g$ leaves $X\supseteq X_\alpha$ from $v\in X_\alpha$, and is
$\alpha$-colored by (2); hence it is the unique identity
$f^0_\alpha$.  If $P$ departed from $X_\alpha$ earlier, its first
departure would again have color $\alpha$, would be distinct from $g$,
and would contradict this uniqueness.  Therefore
\[
 y_0Pv\subseteq X_\alpha\subseteq X.                         \tag{5}
\]

It remains to suppose that $v\notin X_0$.  Unfolding the indexed union in
\noderef{ScheideEnlargedVertexSet}, choose a defective color $\delta$
such that
\[
 v\in U_\delta=\{(R\delta).u1,(R\delta).u2\}.                \tag{6}
\]
Put $i=D.\gammaIndex(\delta)$, $\gamma=C.\gamma(i)$,
$Z=X_0\setminus U_\gamma$, $T_\delta=(R\delta).replacementTree$, and
\[
 \varphi_\delta=\noderef{TwoColorComponentSwap}
   (G,\varphi,\gamma,\delta,\operatorname{inl}v^0_\delta).
\]
The appropriate one of $f^1_\delta,f^2_\delta$ is a
$\delta$-colored identity incident with $v$ and with its other end in
$Z\subseteq X$.  It is distinct from $g$, whose other end is outside
$X$.  The second conclusion of
\noderef{ScheideSymmetricDefectiveInternal} supplies a
$\gamma$-colored identity internal to $U_\delta$; by (6) and
looplessness it too is incident with $v$, and it too is distinct from
$g$.  Properness, (2), and \noderef{AdjacentEdges} therefore give
\[
 \alpha\ne\delta,\qquad \alpha\ne\gamma.                   \tag{7}
\]

We also need both inequalities for $\alpha_0$.  The setup gives
$\alpha_0\ne\gamma$ because $\gamma$ is one of the indexed block
colors.  Since $\delta$ is defective, the definition
\noderef{DefectiveColors} supplies a $\delta$-colored boundary identity
of $X_0$.  Original closedness and (1) forbid such an identity from
having color $\alpha_0$.  Therefore
\[
 \alpha_0\ne\gamma,\qquad \alpha_0\ne\delta.                \tag{8}
\]

We next establish the replacement-tree facts from the nodes which state
them.  The theorem \noderef{ScheideReplacementMaximum} says that
$(T_\delta,e_0,\varphi_\delta)$ is globally maximum.  Applying
\noderef{ScheideGlobalMaximumOrder} to this switched triple gives
\[
 V(T_\delta)\text{ is elementary and closed under }\varphi_\delta.
                                                                    \tag{9}
\]
These conclusions do not come merely from the replacement-data or
one-pair-union definitions.

The exact base and block positions in the structural configuration,
together with vertex injectivity in \noderef{TreeSequence}, show that
$y_0$ is different from the two deleted block vertices.  Hence
$y_0\in Z$.  The setup supplies a full
$(\alpha_0,\delta)$ path from $y_0\in X_0$ to an outside vertex, with
$v^0_\delta$ as its last vertex in $X_0$; it also supplies the missing
color $\alpha_0$ at $y_0$ and the defectiveness of $\delta$.  Unfolding
\noderef{ScheideExitVertex}, these are precisely an exit witness for
$v^0_\delta$, so it belongs to \noderef{ScheideExitVertices}.
Consequently \noderef{ScheideStructuralExitsInBase} puts it in the base,
and the same position argument gives $v^0_\delta\in Z$.  The exact
vertex-set clause of \noderef{ScheideReplacementData} is
\[
 V(T_\delta)=Z\cup U_\delta.                                 \tag{10}
\]
Thus $y_0,v^0_\delta,v\in V(T_\delta)$.

The missing-set formulas in \noderef{ScheideDefectiveSwitchEffect} show
that $\alpha_0$ remains missing at $y_0$.  Indeed, if
$y_0\ne v^0_\delta$, the missing set at $y_0$ is fixed.  If
$y_0=v^0_\delta$, the new missing set is
\[
 \bigl(\overline\varphi(y_0)\setminus\{\gamma\}\bigr)
       \cup\{\delta\};
\]
(1) and $\alpha_0\ne\gamma$ from (8) retain $\alpha_0$ in this set.
Consequently
\[
 \alpha_0\in\noderef{MissingColors}(G,\varphi_\delta,y_0).  \tag{11}
\]

For every edge identity $q$, the defining formula in
\noderef{TwoColorComponentSwap} changes only the two values
$\gamma,\delta$.  Equations (7)--(8) therefore give the two exact
equivalences
\[
 \varphi_\delta(q)=\alpha_0\iff\varphi(q)=\alpha_0,
 \qquad
 \varphi_\delta(q)=\alpha\iff\varphi(q)=\alpha.              \tag{12}
\]
By \noderef{TwoColorIncidenceGraph}, (12) identifies the entire labeled
$(\alpha_0,\alpha)$ incidence graph under the two colorings, incidence by
incidence.  Hence \noderef{TwoColorComponent} identifies the full
component through $y_0$, and the same simple full path $P$ is valid under
$\varphi_\delta$.  Since (10) gives $V(T_\delta)\subseteq X$, while
$v\in V(T_\delta)$ and no vertex of $X$ occurs on $P$ after $v$, the
vertex $v$ is also the last vertex of $T_\delta$ on this switched path.

We now prove the missing-witness assertion needed at $v$.  Let $c$ be any
member of \noderef{TreeUsedColors}$(\varphi_\delta,T_\delta)$.  The exact
edge-set clause in \noderef{ScheideReplacementData} divides the edge
identity witnessing $c$ into a retained old tree edge or one of the two
appended identities $f^1_\delta,f^2_\delta$.  The deleted block begins
at position $p-1\geq1$, so the old first edge remains the first edge of
$T_\delta$ and cannot witness membership in \noderef{TreeUsedColors}.
Moreover, every base vertex belongs to $Z$: the exact base and block
positions are disjoint, by vertex injectivity, just as in the argument
for $y_0$ above.

For a retained colored base edge, the
\noderef{TashkinovTree} field of the original globally maximum triple
supplies an earlier missing-color witness.  Because the edge is still in
the base part, that earlier vertex is a base vertex.  The avoidance clause
of \noderef{ScheideStructuralConfiguration} says that its color differs
from every block color, in particular from the deleted color $\gamma$.
For a retained block edge, the structural configuration says its color is
the corresponding indexed block color and supplies a missing witness in
the base.  Injectivity of the indexed block colors, and the fact that this
is not the deleted block, again show that this retained color differs from
$\gamma$.  In either retained case the edge is incident with $X_0$ and is
not $f^0_\delta$ (every displayed old tree edge has both ends in $X_0$ by
\noderef{TreeSequence}, whereas $f^0_\delta$ has an end outside $X_0$).
Thus \noderef{ScheideDefectiveSwitchEffect} fixes its color.  The same
node fixes the missing set at its base witness unless that witness is
$v^0_\delta$; at $v^0_\delta$, its exact erase-insert formula preserves
the witness because the retained color is not $\gamma$.  Hence every
retained used color is missing under $\varphi_\delta$ at a vertex of $Z$.

For an appended identity, \noderef{ScheideReplacementData} gives its old
color $\delta$, its endpoint in $Z\subseteq X_0$, and its inequality from
$f^0_\delta$.  The fixed-edge clause of
\noderef{ScheideDefectiveSwitchEffect} therefore leaves its color equal
to $\delta$.  The exact missing-set formula at $v^0_\delta$ inserts
$\delta$, and we proved above that $v^0_\delta\in Z$.  Thus the appended
edges also have a missing witness in $Z$.  We have proved
\[
 c\in\noderef{TreeUsedColors}(\varphi_\delta,T_\delta)
 \quad\Longrightarrow\quad
 \exists z\in Z,\ c\in\noderef{MissingColors}(G,\varphi_\delta,z).
                                                                    \tag{13}
\]

Apply \noderef{ScheideGlobalMaximumMissingCounts} to the switched globally
maximum triple and the vertex $v\in V(T_\delta)$.  If $v$ is incident
with $e_0$ its missing set has at least two members; otherwise it has at
least one.  In either case choose an actual color
$\beta\in\noderef{MissingColors}(G,\varphi_\delta,v)$.  If $\beta$ were
used on $T_\delta$, (13) would give $z\in Z$ also missing $\beta$.  But
$z\in X_0$ while $v\notin X_0$, so $z\ne v$; (9), (10), and
\noderef{ColoredVertexSetElementary} forbid two such distinct replacement
tree vertices.  Therefore $\beta$ is unused.  Since it is missing at the
tree vertex $v$, the definitions \noderef{TreeMissingColors} and
\noderef{TreeFreeColors} give
\[
 \beta\in\noderef{TreeFreeColors}(\varphi_\delta,T_\delta).  \tag{14}
\]

Suppose for contradiction that $\alpha$ were defective for
$T_\delta$ under $\varphi_\delta$.  The vertex $y_0$ lies in the tree and
misses $\alpha_0$ by (10)--(11); the path $P$ is the full
$(\alpha_0,\alpha)$ component under $\varphi_\delta$; and $v$ is its last
replacement-tree vertex.  These data satisfy the predicate
\noderef{ScheideExitVertex}, with defective color $\alpha$ and missing
color $\alpha_0$, so $v$ belongs to
\noderef{ScheideExitVertices}$(T_\delta,\varphi_\delta)$.  But
\noderef{ScheideExitNoFree}, applied to the globally maximum replacement
triple, says that such an exit misses no free color.  This contradicts
(14) and the choice that $\beta$ is missing at $v$.  Hence $\alpha$ is
not defective for the replacement tree.

Finally, (9) and (11) forbid every $\alpha_0$-colored boundary identity
of $V(T_\delta)$.  The edge $g$ has one end $v$ in $V(T_\delta)$ and its
other end outside $X\supseteq V(T_\delta)$, so it is an
$\alpha$-colored boundary identity of the replacement tree.  If $P$
left $V(T_\delta)$ before reaching $v$, its first departure would also
have color $\alpha$, would be distinct from $g$, and would make $\alpha$
defective there.  Therefore
\[
 y_0Pv\subseteq V(T_\delta)\subseteq X.                      \tag{15}
\]

Equations (3), (5), and (15) prove in all three cases that $P$ stays in
$X$ through $v$.  By the choice of $v$, no later vertex-node of $P$
belongs to $X$.  Hence $g$ is the only path edge with one endpoint in
$X$ and one outside.  Since $P$ is the entire two-color component and
its edge-nodes retain individual edge identities, an edge-node of that
component belongs to $\partial_G(X)$ exactly when its two neighboring
vertex-nodes form such a crossing.  Thus $g$ is the unique boundary edge
identity asserted.
\end{proof}
